%% Согласно ГОСТ Р 7.0.11-2011:
%% 5.3.3 В заключении диссертации излагают итоги выполненного исследования, рекомендации, перспективы дальнейшей разработки темы.
%% 9.2.3 В заключении автореферата диссертации излагают итоги данного исследования, рекомендации и перспективы дальнейшей разработки темы.
\begin{enumerate}
  \item Проведён анализ производственного процесса целлюлозно-бумажного комбината: описана технологическая цепочка <<БДМ~--- склад полуфабрикатов~--- ПРС/БРС>>, выделены потери материала при~раскрое, переналадки ножей, сроки заказов, многослойность и~двухэтапная резка малых форматов через БРС; сформулированы требования к~производственному плану.

  \item Разработана система дискретно-событийного имитационного моделирования исполнения расписания, учитывающая случайные длительности, отказы и~ремонты машин, обнаруживающая блокировки и~агрегирующая результаты серии прогонов. В~текущей реализации случайные отклонения длительности обработки имеют равномерное распределение, параметры которого задают диапазон отклонений. Выделены показатели времени завершения, запаздывания, просроченных заказов, долей прогонов без опозданий и~с блокировкой, загрузки оборудования и~складов полуфабрикатов.

  \item Предложен порядок сравнения планов в~оптимистичном, реалистичном и~пессимистичном сценариях с~оценкой средних и~квантилей показателей внутри каждого сценария.

  \item Сформулирована математическая модель, в~которой решение представлено трёхкомпонентным глобальным планом $G=(\overline{\overline{X}},\overline{\overline{Y}},\overline{\overline{Z}})$, объединяющим план производства тамбуров на~БДМ, план нарезки на~ПРС/БРС и~расписание операций. Введены параметры заказов, оборудования и~технологические константы.

  \item Определена векторная целевая функция $\mu(\overline{\overline{Y}})=\langle\omega(\overline{\overline{Y}}),\,\varphi(\overline{\overline{Y}}),\,\tau(\overline{\overline{Y}})\rangle\to\min$ и~обоснован лексикографический порядок критериев. Сформулирована система из~десяти ограничений, учитывающая технологическую очередность марок, раскрой, расписание, сроки, наличие полуфабрикатов, содержимое и~зарезервированное содержимое складов полуфабрикатов, а~также совместное размещение оборудования.

  \item Показано, что задача относится к~классу NP-трудных, поскольку классическая задача линейного раскроя является её частным случаем. Обоснована необходимость приближённых численных методов для~промышленных экземпляров.

  \item Установлена научная новизна модели, связанная с~одновременным учётом нескольких ПРС и~БРС, многослойных заказов, двухэтапной нарезки малых форматов, динамики складских запасов отдельных тамбуров, согласованного расписания и~совместного размещения оборудования.

  \item Предложены вспомогательные алгоритмы инкрементального формирования плана нарезки, формирования плана производства и~построения расписания. Они обеспечивают согласованное формирование трёх компонент глобального плана и~сохранение проверяемых ограничений.

  \item Предложен алгоритм формирования случайного допустимого глобального плана, последовательно строящий план нарезки, план производства и~расписание с~сохранением допустимости ограничений модели.

  \item Получены оценки вычислительной сложности этапов численных методов и~генетического поиска. Сложность поиска линейна по~размеру начальной популяции и~суммарному числу рулонов при~фиксированных параметрах оборудования.

  \item Разработан программный комплекс на~языке Python, реализующий полный перебор и~генетический поиск, имитационное моделирование, ввод и~вывод данных в~форматах CSV и~Excel, а~также командный и~графический интерфейсы.

  \item Проведено тестирование программного комплекса, включающее модульные и~интеграционные тесты, экспериментальную валидацию и~проверку свойств имитационного моделирования. Покрытие кода тестами составляет $81\%$. Программный комплекс зарегистрирован в~Реестре программ для~ЭВМ.

  \item Проведены вычислительные эксперименты на~четырёх тестовых экземплярах возрастающей размерности. При~размере популяции $P=10\,000$ генетический поиск воспроизвёл значения целевой функции полного перебора; на~наибольшем экземпляре ускорение составило примерно $59$~раз.

  \item Проведён вычислительный эксперимент на~производственных данных Кондопожского ЦБК: $45$ типов заказов и~$532$ рулона. Построен допустимый план со~значением целевой функции $\langle133525,\,160,\,518\rangle$ и~продолжительностью $86$~ч~$10$~мин; выполнены $10$ имитационных прогонов, средняя фактическая продолжительность составила $3$~сут~$17$~ч~$40$~мин, стандартное отклонение~--- $1$~ч~$07$~мин.

  \item Исследовано влияние размера начальной популяции на~качество решений генетического поиска. При~увеличении $P$ от~$5$ до~$100$ отходы снизились со~$124$ до~$65$, а~дальнейшее увеличение $P$ до~$10\,000$ не~изменило результат, но~увеличило время расчёта.

  \item Результаты внедрены в~практическую деятельность ООО~<<Опти-Софт>>. Предложенные численные методы рекомендованы для~планирования производства и~раскроя на~предприятиях целлюлозно-бумажной промышленности и~в~смежных областях.

  \item Определены ограничения текущей реализации и~направления развития: стохастические факторы учитываются на~этапе имитационной оценки, но~не~при~построении плана; не~рассматриваются межпериодное планирование, интеграция с~корпоративными информационными системами и~параллельное выполнение вычислений.
\end{enumerate}